\documentclass[11pt]{article}
\usepackage{amsmath, amssymb, geometry, graphicx, xcolor}
\geometry{a4paper, margin=1in}

\title{Homework 1.8 Introduction to Linear Transformations}
\author{Tashfeen Omran}
\date{September 2026}

\begin{document}

\maketitle
\tableofcontents
\newpage

% ---------------------------------------------------------
% PART 1: Comprehensive Introduction, Context, and Prerequisites
% ---------------------------------------------------------
\section{Comprehensive Introduction, Context, and Prerequisites}

\subsection{Core Concepts}
The topic of \textbf{Linear Transformations} represents a conceptual shift in linear algebra. Up until now, a matrix equation like $A\mathbf{x} = \mathbf{b}$ has been viewed primarily as a static system of linear equations. In this chapter, we re-imagine the matrix $A$ as an active \textit{function} or \textit{operator} that takes an input vector $\mathbf{x}$ and ``transforms'' it into an output vector $\mathbf{b}$. 

When we define a mapping $T(\mathbf{x}) = A\mathbf{x}$, we introduce the following terminology:
\begin{itemize}
    \item \textbf{Domain:} The set of all possible inputs. If $A$ is an $m \times n$ matrix, it must multiply a vector with $n$ entries. Thus, the domain is $\mathbb{R}^n$.
    \item \textbf{Codomain:} The set of all possible outputs. The result of $A\mathbf{x}$ is a vector with $m$ entries, so the codomain is $\mathbb{R}^m$. We write $T: \mathbb{R}^n \to \mathbb{R}^m$.
    \item \textbf{Image:} The specific output vector $T(\mathbf{x})$ that results from plugging in a specific input $\mathbf{x}$.
    \item \textbf{Range:} The set of \textit{all} actual images. The range is a subset of the codomain.
\end{itemize}

\subsection{Intuition and Derivation}
Think of a linear transformation as a machine that morphs space. When you multiply a vector by a matrix, the matrix can stretch, rotate, shear, or project that vector. 
To be classified as a \textit{linear} transformation, the mapping must preserve the underlying structure of vector addition and scalar multiplication. Geometrically, this means grid lines remain parallel and evenly spaced, and the origin must remain fixed at the origin. If you scale an input by a factor of 2, its output must also scale by exactly a factor of 2. 

\subsection{Historical Context and Motivation}
Historically, mathematicians in the 19th century, such as Arthur Cayley and James Joseph Sylvester, developed matrix algebra to streamline the manipulation of coordinate transformations. They were motivated by problems in projective geometry and celestial mechanics, where keeping track of changing coordinate systems was a monumental task. By viewing matrices as unified transformation operators, they created a powerful language that allowed complex, multi-dimensional changes to be expressed as simple algebraic multiplications. This abstraction is the foundation of modern computer graphics, quantum mechanics, and multivariable calculus.

\subsection{Key Formulas and Properties}
A transformation $T$ is \textbf{linear} if and only if it satisfies both of the following properties for all vectors $\mathbf{u}, \mathbf{v}$ in the domain and all scalars $c$:
\begin{enumerate}
    \item \textbf{Preservation of Addition:} \[ T(\mathbf{u} + \mathbf{v}) = T(\mathbf{u}) + T(\mathbf{v}) \]
    \item \textbf{Preservation of Scalar Multiplication:} \[ T(c\mathbf{u}) = cT(\mathbf{u}) \]
\end{enumerate}
These two properties can be combined into the \textbf{Superposition Principle}:
\[ T(c_1\mathbf{v}_1 + c_2\mathbf{v}_2 + \dots + c_p\mathbf{v}_p) = c_1T(\mathbf{v}_1) + c_2T(\mathbf{v}_2) + \dots + c_pT(\mathbf{v}_p) \]
A critical necessary (though not sufficient) condition for linearity is:
\[ T(\mathbf{0}) = \mathbf{0} \]

\subsection{Prerequisites}
To master this topic, you must be comfortable with:
\begin{itemize}
    \item \textbf{Matrix-Vector Multiplication:} Computing $A\mathbf{x}$ efficiently.
    \item \textbf{Row Reduction (RREF):} Using Gaussian elimination on an augmented matrix $[A \ | \ \mathbf{b}]$ to solve for $\mathbf{x}$.
    \item \textbf{Vector Equations:} Understanding that $A\mathbf{x}$ is a linear combination of the columns of $A$.
    \item \textbf{Parametric Vector Form:} Expressing infinite solution sets using free variables.
\end{itemize}

% ---------------------------------------------------------
% PART 2: Detailed Homework Solutions
% ---------------------------------------------------------
\section{Detailed Homework Solutions}
\textit{Note: As requested, to maximize learning efficiency, we have selected a representative sample covering one or two of every distinct problem type from the exercise set, ensuring all concepts are mastered without repetitive redundancy.}

\subsection{Problem 1 (Calculating Images)}
\textbf{Question:} Let $A = \begin{bmatrix} 2 & 0 \\ 0 & 2 \end{bmatrix}$, and define $T: \mathbb{R}^2 \to \mathbb{R}^2$ by $T(\mathbf{x}) = A\mathbf{x}$. Find the images under $T$ of $\mathbf{u} = \begin{bmatrix} 1 \\ -3 \end{bmatrix}$ and $\mathbf{v} = \begin{bmatrix} a \\ b \end{bmatrix}$.

\textbf{Solution:}
To find the image of a vector, we perform matrix-vector multiplication.
\begin{align*}
    T(\mathbf{u}) &= A\mathbf{u} \\
    &= \begin{bmatrix} 2 & 0 \\ 0 & 2 \end{bmatrix} \begin{bmatrix} 1 \\ -3 \end{bmatrix} \\
    &= \begin{bmatrix} (2)(1) + (0)(-3) \\ (0)(1) + (2)(-3) \end{bmatrix} \\
    &= \begin{bmatrix} 2 \\ -6 \end{bmatrix}
\end{align*}
Now, for the abstract vector $\mathbf{v}$:
\begin{align*}
    T(\mathbf{v}) &= A\mathbf{v} \\
    &= \begin{bmatrix} 2 & 0 \\ 0 & 2 \end{bmatrix} \begin{bmatrix} a \\ b \end{bmatrix} \\
    &= \begin{bmatrix} 2a + 0b \\ 0a + 2b \end{bmatrix} \\
    &= \begin{bmatrix} 2a \\ 2b \end{bmatrix}
\end{align*}
\textbf{Final Answer:} The image of $\mathbf{u}$ is $\begin{bmatrix} 2 \\ -6 \end{bmatrix}$, and the image of $\mathbf{v}$ is $\begin{bmatrix} 2a \\ 2b \end{bmatrix}$. (Geometrically, this represents a scalar dilation by a factor of 2).

\subsection{Problem 3 (Finding Pre-Images and Uniqueness)}
\textbf{Question:} Let $A = \begin{bmatrix} 1 & 0 & -2 \\ -2 & 1 & 6 \\ 3 & -2 & -5 \end{bmatrix}$, $\mathbf{b} = \begin{bmatrix} -1 \\ 7 \\ -3 \end{bmatrix}$. Find a vector $\mathbf{x}$ whose image under $T$ is $\mathbf{b}$, and determine whether $\mathbf{x}$ is unique.

\textbf{Solution:}
We need to solve the equation $T(\mathbf{x}) = \mathbf{b}$, which is equivalent to solving $A\mathbf{x} = \mathbf{b}$. We set up the augmented matrix and row reduce:
\[ \left[ \begin{array}{ccc|c} 1 & 0 & -2 & -1 \\ -2 & 1 & 6 & 7 \\ 3 & -2 & -5 & -3 \end{array} \right] \]
\textit{Step 1:} Eliminate entries below the first pivot (Row 1).
\[ R_2 \leftarrow R_2 + 2R_1 \quad \implies \quad \left[ \begin{array}{ccc|c} 1 & 0 & -2 & -1 \\ 0 & 1 & 2 & 5 \\ 3 & -2 & -5 & -3 \end{array} \right] \]
\[ R_3 \leftarrow R_3 - 3R_1 \quad \implies \quad \left[ \begin{array}{ccc|c} 1 & 0 & -2 & -1 \\ 0 & 1 & 2 & 5 \\ 0 & -2 & 1 & 0 \end{array} \right] \]
\textit{Step 2:} Eliminate entry below the second pivot (Row 2).
\[ R_3 \leftarrow R_3 + 2R_2 \quad \implies \quad \left[ \begin{array}{ccc|c} 1 & 0 & -2 & -1 \\ 0 & 1 & 2 & 5 \\ 0 & 0 & 5 & 10 \end{array} \right] \]
\textit{Step 3:} Solve the resulting system (backward substitution).
\begin{align*}
    5x_3 &= 10 \implies x_3 = 2 \\
    x_2 + 2(2) &= 5 \implies x_2 = 1 \\
    x_1 - 2(2) &= -1 \implies x_1 = 3
\end{align*}
The vector is $\mathbf{x} = \begin{bmatrix} 3 \\ 1 \\ 2 \end{bmatrix}$. 
Because there is a pivot in every column of the coefficient matrix, there are no free variables. Therefore, the solution is unique.

\textbf{Final Answer:} $\mathbf{x} = \begin{bmatrix} 3 \\ 1 \\ 2 \end{bmatrix}$, and it is unique.

\subsection{Problem 7 (Determining Dimensions)}
\textbf{Question:} Let $A$ be a $4 \times 6$ matrix. What must $a$ and $b$ be in order to define $T: \mathbb{R}^a \to \mathbb{R}^b$ by $T(\mathbf{x}) = A\mathbf{x}$?

\textbf{Solution:}
The matrix $A$ has dimensions $4 \times 6$ (4 rows and 6 columns). 
For the matrix-vector product $A\mathbf{x}$ to be defined, the vector $\mathbf{x}$ must have exactly as many entries as $A$ has columns. Therefore, $\mathbf{x}$ must be a vector in $\mathbb{R}^6$, meaning the domain is $\mathbb{R}^6$. Thus, $a = 6$.
The output of $A\mathbf{x}$ is a linear combination of the columns of $A$. Since each column has 4 entries (corresponding to the 4 rows), the output vector exists in $\mathbb{R}^4$. The codomain is $\mathbb{R}^4$. Thus, $b = 4$.

\textbf{Final Answer:} $a = 6$ and $b = 4$.

\subsection{Problem 9 (Mapping to the Zero Vector)}
\textbf{Question:} Let $A = \begin{bmatrix} 1 & -4 & 7 & -5 \\ 0 & 1 & -4 & 3 \\ 2 & -6 & 6 & -4 \end{bmatrix}$. Find all $\mathbf{x}$ in $\mathbb{R}^4$ that are mapped into the zero vector by the transformation $\mathbf{x} \mapsto A\mathbf{x}$.

\textbf{Solution:}
We must solve the homogeneous system $A\mathbf{x} = \mathbf{0}$. We augment $A$ with the zero vector and row reduce:
\[ \left[ \begin{array}{cccc|c} 1 & -4 & 7 & -5 & 0 \\ 0 & 1 & -4 & 3 & 0 \\ 2 & -6 & 6 & -4 & 0 \end{array} \right] \]
\textit{Step 1:} Eliminate below the first pivot.
\[ R_3 \leftarrow R_3 - 2R_1 \quad \implies \quad \left[ \begin{array}{cccc|c} 1 & -4 & 7 & -5 & 0 \\ 0 & 1 & -4 & 3 & 0 \\ 0 & 2 & -8 & 6 & 0 \end{array} \right] \]
\textit{Step 2:} Eliminate below the second pivot.
\[ R_3 \leftarrow R_3 - 2R_2 \quad \implies \quad \left[ \begin{array}{cccc|c} 1 & -4 & 7 & -5 & 0 \\ 0 & 1 & -4 & 3 & 0 \\ 0 & 0 & 0 & 0 & 0 \end{array} \right] \]
\textit{Step 3:} Move to Reduced Row Echelon Form (RREF) to read off the parametric solution.
\[ R_1 \leftarrow R_1 + 4R_2 \quad \implies \quad \left[ \begin{array}{cccc|c} 1 & 0 & -9 & 7 & 0 \\ 0 & 1 & -4 & 3 & 0 \\ 0 & 0 & 0 & 0 & 0 \end{array} \right] \]
The variables $x_1$ and $x_2$ are basic (pivot columns 1 and 2). Variables $x_3$ and $x_4$ are free.
\begin{align*}
    x_1 - 9x_3 + 7x_4 &= 0 \implies x_1 = 9x_3 - 7x_4 \\
    x_2 - 4x_3 + 3x_4 &= 0 \implies x_2 = 4x_3 - 3x_4 \\
    x_3 &= x_3 \quad \text{(free)} \\
    x_4 &= x_4 \quad \text{(free)}
\end{align*}
Expressing this in parametric vector form:
\[ \mathbf{x} = \begin{bmatrix} 9x_3 - 7x_4 \\ 4x_3 - 3x_4 \\ x_3 \\ x_4 \end{bmatrix} = x_3\begin{bmatrix} 9 \\ 4 \\ 1 \\ 0 \end{bmatrix} + x_4\begin{bmatrix} -7 \\ -3 \\ 0 \\ 1 \end{bmatrix} \]

\textbf{Final Answer:} $\mathbf{x} = x_3\begin{bmatrix} 9 \\ 4 \\ 1 \\ 0 \end{bmatrix} + x_4\begin{bmatrix} -7 \\ -3 \\ 0 \\ 1 \end{bmatrix}$, for any free parameters $x_3, x_4 \in \mathbb{R}$.

\subsection{Problem 17 (Linearity Properties)}
\textbf{Question:} Let $T: \mathbb{R}^2 \to \mathbb{R}^2$ be a linear transformation that maps $\mathbf{u} = \begin{bmatrix} 2 \\ 1 \end{bmatrix}$ into $\begin{bmatrix} 1 \\ 3 \end{bmatrix}$ and maps $\mathbf{v} = \begin{bmatrix} -1 \\ 2 \end{bmatrix}$ into $\begin{bmatrix} 1 \\ -5 \end{bmatrix}$. Use the fact that $T$ is linear to find the images of $5\mathbf{u}$, $4\mathbf{v}$, and $5\mathbf{u} + 4\mathbf{v}$.

\textbf{Solution:}
We are given $T(\mathbf{u}) = \begin{bmatrix} 1 \\ 3 \end{bmatrix}$ and $T(\mathbf{v}) = \begin{bmatrix} 1 \\ -5 \end{bmatrix}$.
Because $T$ is linear, we can use the property $T(c\mathbf{u}) = cT(\mathbf{u})$.
\begin{align*}
    T(5\mathbf{u}) &= 5T(\mathbf{u}) = 5\begin{bmatrix} 1 \\ 3 \end{bmatrix} = \begin{bmatrix} 5 \\ 15 \end{bmatrix}
\end{align*}
\begin{align*}
    T(4\mathbf{v}) &= 4T(\mathbf{v}) = 4\begin{bmatrix} 1 \\ -5 \end{bmatrix} = \begin{bmatrix} 4 \\ -20 \end{bmatrix}
\end{align*}
Because $T$ is linear, we can also use the property $T(\mathbf{a} + \mathbf{b}) = T(\mathbf{a}) + T(\mathbf{b})$.
\begin{align*}
    T(5\mathbf{u} + 4\mathbf{v}) &= T(5\mathbf{u}) + T(4\mathbf{v}) \\
    &= \begin{bmatrix} 5 \\ 15 \end{bmatrix} + \begin{bmatrix} 4 \\ -20 \end{bmatrix} \\
    &= \begin{bmatrix} 9 \\ -5 \end{bmatrix}
\end{align*}

\textbf{Final Answer:} $T(5\mathbf{u}) = \begin{bmatrix} 5 \\ 15 \end{bmatrix}$, $T(4\mathbf{v}) = \begin{bmatrix} 4 \\ -20 \end{bmatrix}$, and $T(5\mathbf{u} + 4\mathbf{v}) = \begin{bmatrix} 9 \\ -5 \end{bmatrix}$.

\subsection{Problem 40 (Non-Linear Transformations)}
\textbf{Question:} Show that the transformation $T$ defined by $T(x_1, x_2) = (4x_1 - 2x_2, 3|x_2|)$ is not linear.

\textbf{Solution:}
To prove a transformation is not linear, we only need to find a single counterexample that violates one of the linearity properties. The absolute value function is suspicious because it does not preserve scalar multiplication for negative scalars. 
Let us test the property $T(c\mathbf{u}) = cT(\mathbf{u})$.
Let $\mathbf{u} = (0, 1)$ and $c = -1$.
First, compute the left side, $T(c\mathbf{u})$:
\[ c\mathbf{u} = -1(0, 1) = (0, -1) \]
\[ T(0, -1) = (4(0) - 2(-1), 3|-1|) = (2, 3(1)) = (2, 3) \]
Now, compute the right side, $cT(\mathbf{u})$:
\[ T(\mathbf{u}) = T(0, 1) = (4(0) - 2(1), 3|1|) = (-2, 3) \]
\[ cT(\mathbf{u}) = -1(-2, 3) = (2, -3) \]
Since $T(-\mathbf{u}) = (2, 3)$ and $-T(\mathbf{u}) = (2, -3)$, we have $T(-\mathbf{u}) \neq -T(\mathbf{u})$. The scalar multiplication property is violated.

\textbf{Final Answer:} $T$ is not linear because it fails the scalar multiplication property $T(c\mathbf{u}) = cT(\mathbf{u})$ when $c < 0$, as demonstrated by the counterexample $T(0, -1) = (2, 3) \neq -( -2, 3 ) = -T(0, 1)$.

% ---------------------------------------------------------
% PART 3: In-Depth Analysis of Problems and Techniques
% ---------------------------------------------------------
\section{In-Depth Analysis of Problems and Techniques}

\subsection{A) Problem Types and General Approach}
\begin{itemize}
    \item \textbf{Type 1: Forward Mapping (Image Calculation):} (e.g., Prob 1) You are given $A$ and $\mathbf{x}$. Strategy: Simply compute the matrix-vector product $A\mathbf{x}$.
    \item \textbf{Type 2: Inverse Mapping (Finding Pre-images):} (e.g., Prob 3, 9) You are given $A$ and the output $\mathbf{b}$ (or $\mathbf{0}$). Strategy: Set up the augmented matrix $[A \ | \ \mathbf{b}]$, row reduce to echelon form. If consistent, use back-substitution or parametric vector form to find the inputs $\mathbf{x}$.
    \item \textbf{Type 3: Domain and Codomain Identification:} (e.g., Prob 7) Analyzing the dimensions of matrix $A$. Strategy: Remember the rule ``$m \times n$ matrix maps $\mathbb{R}^n$ to $\mathbb{R}^m$''. The number of columns is the domain dimension; the number of rows is the codomain dimension.
    \item \textbf{Type 4: Abstract Property Application:} (e.g., Prob 17, 40) Working directly with $T(\mathbf{u}+\mathbf{v})$ and $T(c\mathbf{u})$ without using a matrix. Strategy: Distribute $T$ across linear combinations or search for a counterexample (negative numbers usually break non-linear functions like absolute values or squares).
\end{itemize}

\subsection{B) Key Algebraic and Calculus Manipulations}
\begin{itemize}
    \item \textbf{Gaussian Elimination and RREF:} Essential for Problems 3 and 9. The manipulation turns a complex system of equations into an upper triangular form, revealing whether a solution exists, is unique, or contains free variables.
    \item \textbf{Parametric Vector Decomposition:} In Problem 9, after finding $x_1 = 9x_3 - 7x_4$ and $x_2 = 4x_3 - 3x_4$, the crucial algebraic step is factoring out the free variables $x_3$ and $x_4$ as scalar weights on basis vectors. This manipulation is foundational for finding the null space (kernel) of a transformation.
    \item \textbf{Counter-Example Construction:} In Problem 40, the trick was deliberately choosing a negative scalar. The absolute value function $|x_2|$ masks the sign of the input. Feeding it a negative scalar exposes the non-linearity.
\end{itemize}

% ---------------------------------------------------------
% PART 4: "Cheatsheet" and Tips for Success
% ---------------------------------------------------------
\section{"Cheatsheet" and Tips for Success}

\subsection{Essential Formulas}
\begin{itemize}
    \item \textbf{Matrix Mapping:} $T(\mathbf{x}) = A\mathbf{x}$
    \item \textbf{Linearity Properties:} $T(c\mathbf{u} + d\mathbf{v}) = cT(\mathbf{u}) + dT(\mathbf{v})$
    \item \textbf{The Origin Rule:} $T(\mathbf{0}) = \mathbf{0}$. (If a function maps the zero vector to anything other than the zero vector, it is immediately NOT linear).
\end{itemize}

\subsection{Tips, Tricks, and Shortcuts}
\begin{itemize}
    \item \textbf{Dimension Mismatch Check:} Always check your matrix dimensions before starting an operation. If $A$ is $m \times n$, then $\mathbf{x}$ MUST be $n \times 1$. If they don't match, you've set up the problem wrong.
    \item \textbf{Fast Uniqueness Test:} When finding a pre-image, look at the columns of the echelon form. If every column of the coefficient matrix has a pivot, the solution is unique. If there is a column without a pivot, there are free variables, meaning infinite solutions.
    \item \textbf{The "Plus a Constant" Trap:} A transformation like $T(\mathbf{x}) = A\mathbf{x} + \mathbf{b}$ (where $\mathbf{b} \neq \mathbf{0}$) is called an \textit{affine} transformation. It is \textbf{not} a linear transformation because $T(\mathbf{0}) = \mathbf{b} \neq \mathbf{0}$. 
\end{itemize}

% ---------------------------------------------------------
% PART 5: Conceptual Synthesis and The "Big Picture"
% ---------------------------------------------------------
\section{Conceptual Synthesis and The "Big Picture"}

\subsection{A) Thematic Connections}
The core theme of this chapter is the transition from viewing algebra as "finding an unknown variable" to viewing it as "studying the behavior of a function mapping spaces". This parallels the shift in calculus from studying single algebraic equations to studying functions $f(x)$. Here, the function is $T(\mathbf{x})$ and the input is a multi-dimensional vector. 

\subsection{B) Forward and Backward Links}
\textbf{Backward:} This chapter acts as the geometrical realization of solving $A\mathbf{x} = \mathbf{b}$. We are interpreting the mechanics of Chapter 1.1-1.5 through the lens of coordinate geometry. \\
\textbf{Forward:} The concept of linear transformations is the absolute bedrock for the rest of linear algebra. In coming chapters, you will study:
\begin{itemize}
    \item \textbf{Eigenvalues/Eigenvectors:} Finding vectors that don't change direction when transformed ($T(\mathbf{x}) = \lambda \mathbf{x}$).
    \item \textbf{Jacobian Matrices:} In multivariable calculus, non-linear mappings are approximated by linear matrix transformations using the Jacobian matrix.
\end{itemize}

% ---------------------------------------------------------
% PART 6: Real-World Application and Modeling
% ---------------------------------------------------------
\section{Real-World Application and Modeling}

\subsection{A) Concrete Scenarios (Finance and Statistical Focus)}
\begin{enumerate}
    \item \textbf{Arbitrage Pricing Theory (APT) and Factor Models:} In quantitative finance, the return of a portfolio is often modeled as a linear combination of underlying macroeconomic risk factors (e.g., inflation, GDP growth, interest rates). A linear transformation maps a vector of factor returns into a vector of specific asset returns. An analyst must solve the inverse problem (finding the pre-image) to determine what portfolio weights will yield a desired target return profile.
    \item \textbf{Credit Rating Markov Chains:} In credit risk modeling, the state of a bond (e.g., AAA, BBB, Default) changes over time. We represent the probabilities of a bond being in any given state as a state vector $\mathbf{x}$. The transition probabilities from one year to the next are represented by a matrix $A$. The mapping $\mathbf{x}_{t+1} = A\mathbf{x}_t$ is a linear transformation that predicts the future credit distribution of a portfolio.
    \item \textbf{Stochastic Calculus - Covariance Transformations:} When pricing derivatives with multiple correlated underlying assets (like a basket option), quants use a linear transformation (specifically the Cholesky decomposition matrix) to transform a vector of independent random standard normal variables into a vector of correlated random variables used in Monte Carlo simulations.
\end{enumerate}

\subsection{B) Model Problem Setup: Credit Rating Markov Chain}
\textbf{The Scenario:} An investment bank holds a portfolio of bonds that can be in one of three states: Investment Grade (IG), Speculative Grade (SG), and Default (D). We want to find the expected distribution of these bonds in one year.
\textbf{Variables:}
Let $\mathbf{x} = \begin{bmatrix} x_1 \\ x_2 \\ x_3 \end{bmatrix}$ be the state vector, where $x_1$ is the proportion of IG bonds, $x_2$ is SG, and $x_3$ is D. \\
Let $A$ be the transition matrix where entry $a_{ij}$ is the probability of moving from state $j$ to state $i$.
For example:
\[ A = \begin{bmatrix} 0.90 & 0.15 & 0.00 \\ 0.08 & 0.75 & 0.00 \\ 0.02 & 0.10 & 1.00 \end{bmatrix} \]
\textbf{The Integral/Equation to Solve:}
If the bank starts with a portfolio entirely in Investment Grade, the initial vector is $\mathbf{x}_0 = \begin{bmatrix} 1 \\ 0 \\ 0 \end{bmatrix}$. To find the distribution after one year, the bank evaluates the linear transformation:
\[ T(\mathbf{x}_0) = A\mathbf{x}_0 = \begin{bmatrix} 0.90 & 0.15 & 0.00 \\ 0.08 & 0.75 & 0.00 \\ 0.02 & 0.10 & 1.00 \end{bmatrix} \begin{bmatrix} 1 \\ 0 \\ 0 \end{bmatrix} \]
To find the state in 10 years, they would compute $T^{10}(\mathbf{x}_0) = A^{10}\mathbf{x}_0$.

% ---------------------------------------------------------
% PART 7: Common Variations and Untested Concepts
% ---------------------------------------------------------
\section{Common Variations and Untested Concepts}

A vital concept introduced in the very next section (1.9), but intrinsically tied to 1.8, is finding the \textbf{Standard Matrix of a Linear Transformation}. Often, tests will describe a transformation geometrically (e.g., "rotates space by 90 degrees") or by its action on basis vectors, and ask you to construct the matrix $A$.

\textbf{Untested Concept:} Constructing $A$ from $T(\mathbf{e}_1)$ and $T(\mathbf{e}_2)$.
\textbf{Explanation:} The columns of the matrix $A$ are simply the images of the standard basis vectors. If you know where the axes go, you know where the entire space goes.

\textbf{Example:} Let $T: \mathbb{R}^2 \to \mathbb{R}^2$ be a linear transformation that maps $\mathbf{e}_1 = \begin{bmatrix} 1 \\ 0 \end{bmatrix}$ to $\begin{bmatrix} 3 \\ 5 \end{bmatrix}$ and maps $\mathbf{e}_2 = \begin{bmatrix} 0 \\ 1 \end{bmatrix}$ to $\begin{bmatrix} -1 \\ 2 \end{bmatrix}$. Find the standard matrix $A$.
\textbf{Solution:} 
The matrix $A$ is formed by concatenating the images as columns:
\[ A = \big[ T(\mathbf{e}_1) \ \ \ T(\mathbf{e}_2) \big] = \begin{bmatrix} 3 & -1 \\ 5 & 2 \end{bmatrix} \]

% ---------------------------------------------------------
% PART 8: Advanced Diagnostic Testing: "Find the Flaw"
% ---------------------------------------------------------
\section{Advanced Diagnostic Testing: "Find the Flaw"}
\textit{For each problem, a flawed solution is provided. Act as the grader: locate the specific error, explain it in one sentence, and provide the correct calculation.}

\subsection{Flawed Problem 1}
\textbf{Question:} Let $T(x, y) = (x + 2y, y + 1)$. Is $T$ a linear transformation? Prove it. \\
\textbf{Flawed Solution:} Yes, $T$ is linear. We can check scalar multiplication. Let $c = 2$ and $\mathbf{u} = (1, 1)$. $T(\mathbf{u}) = (3, 2)$. Multiplying by $c$ gives $2T(\mathbf{u}) = (6, 4)$. Now compute $T(2\mathbf{u}) = T(2, 2) = (2 + 2(2), 2 + 1) = (6, 3)$. Wait, $(6, 4) \neq (6, 3)$. Therefore, the arithmetic was done wrong in the expansion, but it is linear because it's composed of linear equations. \\
\textbf{The Flaw:} The conclusion incorrectly dismisses the failed scalar multiplication test; the transformation is actually non-linear due to the affine $+1$ term. \\
\textbf{The Correction:} $T$ is not linear. The fastest way to prove this is checking the origin: $T(0, 0) = (0 + 0, 0 + 1) = (0, 1)$. Since $T(\mathbf{0}) \neq \mathbf{0}$, $T$ cannot be a linear transformation.

\subsection{Flawed Problem 2}
\textbf{Question:} Let $A = \begin{bmatrix} 1 & 2 & 3 \\ 4 & 5 & 6 \end{bmatrix}$. The transformation $T(\mathbf{x}) = A\mathbf{x}$ maps from $\mathbb{R}^m$ to $\mathbb{R}^n$. Find $m$ and $n$. \\
\textbf{Flawed Solution:} The matrix $A$ has 2 rows and 3 columns. Therefore, it is a $2 \times 3$ matrix. This means it takes a vector with 2 entries and transforms it into a vector with 3 entries. So $T: \mathbb{R}^2 \to \mathbb{R}^3$, making $m = 2$ and $n = 3$. \\
\textbf{The Flaw:} The dimensions for the domain and codomain are reversed because a $2 \times 3$ matrix must multiply a $3 \times 1$ vector. \\
\textbf{The Correction:} The matrix $A$ must multiply a vector $\mathbf{x}$ in $\mathbb{R}^3$ (since $A$ has 3 columns), and the output will be in $\mathbb{R}^2$ (since $A$ has 2 rows). Therefore, $T: \mathbb{R}^3 \to \mathbb{R}^2$, meaning $m = 3$ and $n = 2$.

\subsection{Flawed Problem 3}
\textbf{Question:} Let $T: \mathbb{R}^3 \to \mathbb{R}^3$ be linear. If $T(\mathbf{v}) = \mathbf{0}$ and $\mathbf{v} = \begin{bmatrix} 1 \\ -1 \\ 2 \end{bmatrix}$, find $T(5\mathbf{v})$. \\
\textbf{Flawed Solution:} By linearity, $T(5\mathbf{v}) = 5 + T(\mathbf{v})$. Since $T(\mathbf{v}) = \mathbf{0}$, we have $T(5\mathbf{v}) = 5 + \mathbf{0} = \begin{bmatrix} 5 \\ 5 \\ 5 \end{bmatrix}$. \\
\textbf{The Flaw:} The scalar multiplication property is $T(c\mathbf{v}) = cT(\mathbf{v})$, not addition of the scalar. \\
\textbf{The Correction:} Using the correct linearity property, $T(5\mathbf{v}) = 5T(\mathbf{v}) = 5(\mathbf{0}) = \begin{bmatrix} 0 \\ 0 \\ 0 \end{bmatrix}$.

\subsection{Flawed Problem 4}
\textbf{Question:} Evaluate the matrix transformation $T(\mathbf{x}) = A\mathbf{x}$ where $A = \begin{bmatrix} 2 & 1 \\ 0 & 3 \end{bmatrix}$ and $\mathbf{x} = \begin{bmatrix} 4 \\ 5 \end{bmatrix}$. \\
\textbf{Flawed Solution:} 
\[ A\mathbf{x} = \begin{bmatrix} 2 & 1 \\ 0 & 3 \end{bmatrix} \begin{bmatrix} 4 \\ 5 \end{bmatrix} = \begin{bmatrix} (2)(4) \\ (1)(5) \\ (0)(4) \\ (3)(5) \end{bmatrix} = \begin{bmatrix} 8 \\ 5 \\ 0 \\ 15 \end{bmatrix} \]
\textbf{The Flaw:} The matrix multiplication was performed element-wise by flattening the matrices, instead of taking the dot product of rows with the column vector. \\
\textbf{The Correction:} Multiply the rows of $A$ by the column of $\mathbf{x}$:
\[ A\mathbf{x} = \begin{bmatrix} 2(4) + 1(5) \\ 0(4) + 3(5) \end{bmatrix} = \begin{bmatrix} 8 + 5 \\ 0 + 15 \end{bmatrix} = \begin{bmatrix} 13 \\ 15 \end{bmatrix} \]

\subsection{Flawed Problem 5}
\textbf{Question:} Given $A = \begin{bmatrix} 1 & -1 \\ 2 & -2 \end{bmatrix}$ and $\mathbf{b} = \begin{bmatrix} 3 \\ 6 \end{bmatrix}$. Is there a unique $\mathbf{x}$ such that $T(\mathbf{x}) = \mathbf{b}$? \\
\textbf{Flawed Solution:} Yes, it is unique. I found that $\mathbf{x} = \begin{bmatrix} 4 \\ 1 \end{bmatrix}$ works because $\begin{bmatrix} 1 & -1 \\ 2 & -2 \end{bmatrix} \begin{bmatrix} 4 \\ 1 \end{bmatrix} = \begin{bmatrix} 4-1 \\ 8-2 \end{bmatrix} = \begin{bmatrix} 3 \\ 6 \end{bmatrix}$. Since I found a vector that works, the solution exists and is unique. \\
\textbf{The Flaw:} Finding a single valid vector proves existence but not uniqueness, which requires checking the RREF of the matrix for free variables. \\
\textbf{The Correction:} Set up the augmented matrix $[A \ | \ \mathbf{b}] = \left[ \begin{array}{cc|c} 1 & -1 & 3 \\ 2 & -2 & 6 \end{array} \right]$. Row reducing ($R_2 - 2R_1$) yields $\left[ \begin{array}{cc|c} 1 & -1 & 3 \\ 0 & 0 & 0 \end{array} \right]$. Because the second column lacks a pivot, $x_2$ is a free variable. The system has infinitely many solutions, so the vector $\mathbf{x}$ is \textbf{not unique}. (For example, $\mathbf{x} = \begin{bmatrix} 3 \\ 0 \end{bmatrix}$ also works).

\end{document}